% Q02_solutions.tex — REFERENCE SOLUTIONS to paper/Q02_problems.tex.
% Author-only and read-denied in study mode (Read(./paper/solutions/**)).
% Built only by `make -C paper paper-solutions` (scripts/ci.sh, or the
% author's linked worktree); `make paper` never touches this directory.
%
% Every verifier result quoted below was captured on the course pins (Dafny
% 4.11.0, Verus 0.2026.09.27.3cf1832) through scripts/vcheck, one solver job at
% a time under a 16G cap, by the DV2b kit lane for Q02 on 2026-10-05. The probe
% files are archived under
% ~/.cache/dafny-verus-learn-0/author-probes/dv2b/q02-kit/ (README there).
% B3's fields are graded against the learner's own cheatsheets.md; this file
% records which lines of the pinned range each field draws on and never copies
% the card.
\documentclass[11pt]{article}
\input{preamble}

\title{Q02 --- First syntax I \\ \large Reference solutions (author-only)}
\author{dafny-verus-learn-0}
\date{}

\begin{document}
\maketitle

\section*{How these are graded}
B1 is informal-proof-first (the five proof dimensions; its blueprint is
re-scored by \verb+--post+ against the learner's verified H3d and H2v). B2 is a
concept answer (\texttt{correctness}: the grader tries to counterexample it,
with the rubric's Probe B). B3 is a paper card (\texttt{card-fidelity}, per
field). The solutions below give the answers; the rubric never states them.
The structure of B1 matches the formal references of H3d and H2v (spec
section~4.4's author rule).

% ---------------------------------------------------------------------------
\section*{B1 (informal-proof-first)}

\textbf{(a) Statement.} For every $s$ of type \dfy{seq<nat>},
$\mathit{ToNat}(\mathit{Normalize}(s)) = \mathit{ToNat}(s)$.

\textbf{Induction on} the length $|s|$ (the lemma's \dfy{decreases |s|}),
following the three cases of \dfy{Normalize}.

\textbf{IH (as a formula).} For every $t$ with $|t| < |s|$:
$\mathit{ToNat}(\mathit{Normalize}(t)) = \mathit{ToNat}(t)$. It is used once,
at $t = p$ below, where $|p| = |s| - 1$.

\textbf{Base and the two non-recursive cases.}
\begin{itemize}
\item $|s| = 0$: $\mathit{Normalize}(s) = [\,] = s$, so both sides are
  $\mathit{ToNat}(s)$ (indeed 0).
\item $|s| > 0$ and $s[|s|-1] \neq 0$: $\mathit{Normalize}(s) = s$ by
  definition; nothing to prove.
\end{itemize}

\textbf{Step} ($|s| > 0$ and $s[|s|-1] = 0$). Let $p = s[..|s|-1]$, the
sequence without its top limb. Then
\begin{enumerate}
\item $\mathit{Normalize}(s) = \mathit{Normalize}(p)$, by the definition's third
  case;
\item $s = p + [0]$, as sequences: the same length, and the same limb at every
  index (limb $|s|-1$ is the zero);
\item $\mathit{ToNat}(p + [0]) = \mathit{ToNat}(p)$, L1's worked lemma
  \dfy{ToNatAppendZero} at $p$; with item~2, $\mathit{ToNat}(s) =
  \mathit{ToNat}(p)$;
\item the IH at $p$: $\mathit{ToNat}(\mathit{Normalize}(p)) = \mathit{ToNat}(p)$.
\end{enumerate}
Chaining 1, 4 and 3: $\mathit{ToNat}(\mathit{Normalize}(s)) =
\mathit{ToNat}(\mathit{Normalize}(p)) = \mathit{ToNat}(p) = \mathit{ToNat}(s)$.

\textbf{Helper lemmas (the blueprint).} (i) L1's \dfy{ToNatAppendZero(p)}:
$\mathit{ToNat}(p + [0]) = \mathit{ToNat}(p)$; (ii) the sequence equation
$s = p + [0]$ for $p = s[..|s|-1]$ when the top limb is zero, stated as an
assertion (the solver relates $\mathit{ToNat}(p + [0])$ to $\mathit{ToNat}(s)$
only once it is told the two sequences are equal). Items (i) and (ii) together
are exactly H1d's statement, $\mathit{ToNat}(s[..|s|-1]) = \mathit{ToNat}(s)$,
so naming H1d as the one helper is equally correct. No other helper.

\textbf{Gaps.} None.

\textbf{The formal reference it matches} (H3d's twin):
\begin{dafnycode}
lemma NormalizeToNat(s: seq<nat>)
  ensures ToNat(Normalize(s)) == ToNat(s)
  decreases |s|
{
  if |s| > 0 && s[|s| - 1] == 0 {
    var p := s[..|s| - 1];
    assert s == p + [0];
    ToNatAppendZero(p);
    NormalizeToNat(p);
  }
}
\end{dafnycode}
The two non-recursive cases need no line: the verifier unfolds
\dfy{Normalize} once.

\textbf{(b) Verus.} The loop is the iterative form of the same induction.
\begin{itemize}
\item \textbf{Invariant:} $\mathit{value}(a) = \mathit{value}(a_{\mathit{old}})$,
  written \vrs{a.to_nat() == old(a).to_nat()}: the value at every pass is the
  value at the call. It plays the IH; each pass is one case-3 step.
\item \textbf{On entry:} \vrs{a} is \vrs{old(a)}.
\item \textbf{Kept by a pass:} the guard says the vector is non-empty with a zero
  top limb; the given call (H1v's statement) gives
  \vrs{seq_to_nat(a@.drop_last()) == seq_to_nat(a@)}, and \vrs{pop} makes the
  new view \vrs{a@.drop_last()}, so the value is unchanged.
\item \textbf{At the exit:} the guard is false, so the vector is empty or its top
  limb is non-zero, which is \vrs{wf()}; the invariant gives the value clause.
  The measure \vrs{a.limbs.len()} drops by one per pass.
\end{itemize}
Verus checks the body from the invariant and the guard alone, so without the
invariant nothing ties \vrs{a} inside the loop to \vrs{old(a)}. The twin's
invariant is that one clause.

% ---------------------------------------------------------------------------
\section*{B2 (concept)}

\textbf{(a)} A clause is inductive when it holds on entry and every pass keeps
it: from any state where the clause (with the other clauses kept beside it)
and the guard hold, the body re-establishes it. So the \emph{on entry} and
\emph{maintained} obligations decide inductiveness; the \emph{at exit}
obligation decides only whether the invariant is strong enough for what
follows the loop (an inductive invariant can still be too weak for the
postcondition). Inductiveness is relative to the set of clauses kept.

\textbf{(b)} One example (any loop with the same shape is correct):
\begin{dafnycode}
k := 0;
t := 0;
while k < m        // m: nat
{
  k := k + 1;
  t := t + 1;
}
\end{dafnycode}
Clause: $t \leq m$. Every reachable state has $t = k \leq m$, so the clause
holds there. Alone it is not inductive: from $k = 0$, $t = 1$, $m = 1$ (the
clause and the guard hold), one pass gives $t = 2 > m$. The rubric's Probe B
runs exactly this pass as a compiled test (captured: \texttt{1 passed, 0
failed}); the same loop with the clauses $k \leq m$, $t = k$ beside $t \leq m$
verifies (Probe B$'$, captured: \texttt{verified}).

\textbf{(c)} Pruning re-checks the clauses it holds: with $t \leq m$ alone the
maintenance check fails, so pruning drops it although it is true at every
reachable state. Beside $k \leq m$ and $t = k$ the set is inductive and pruning
keeps all three. Dropping a true clause is sound but can lose what the
postcondition needs; the card then restores the attempt.

% ---------------------------------------------------------------------------
\section*{B3 (paper card 02.7): reference fields and their lines}

Pinned range: \texttt{cheatsheets.md} lines 766--791 (computed 2026-10-05),
\[ \texttt{sha256 04b270b5bf2545aa6727055642d13097d8abedc941d775dc6dc78ec3ae7be6cc}. \]
A field is correct when it agrees with the lines named; wording is free. Each
entry: the reference content (paraphrased), then the lines it draws on.

\begin{description}
\item[Task shape.] Input: a Dafny program whose specification
  (\dfy{requires}/\dfy{ensures}) is given and whose annotations are missing.
  Output: the program with loop invariants, assertions and \dfy{decreases}
  clauses added, code unchanged, so that Dafny verifies it; an LLM in a
  feedback loop with the verifier. \emph{Lines 770, 772.}
\item[Oracle.] Two checks: the Dafny verifier (SMT-backed) accepts the attempt,
  and the diff-check: stripping every annotation from the attempt, with Dafny's
  own parser, gives back the original program. \emph{Lines 772, 775.}
\item[What is trusted.] Dafny 4.11.0 and its solver; Dafny's parser inside the
  strip; the benchmarks' specifications as the intent (their base programs were
  re-derived, fixing 70 of 782 that did not compile); whether spec edits or
  \dfy{assume} are rejected is not stated, so they are trusted silently.
  \emph{Lines 772, 775, 789.}
\item[Cheat surface.] Closed: changing the code (the card's example rewrites an
  assignment to drop a \texttt{+1}), caught by the diff-check. Open or
  unstated: editing the specification, adding \dfy{assume}; hints that may
  encode benchmark knowledge. \emph{Lines 775, 789.}
\item[Metric.] Percentage of programs verified, per benchmark (Clover, MBPP,
  HumanEval, DafnyBench); the headline 86.18\% on DafnyBench with Claude 3.5;
  ablation gains in percentage points; 69.5\% for the fine-tuned small model on
  197 held-out programs. \emph{Lines 770, 780--787.}
\item[Budget.] $N = 10$ attempts per program for the Claude runs; $N = 5$ for
  the held-out small-model runs; no cost, token or time figure.
  \emph{Lines 772, 787, 789.}
\item[Variance.] One run per configuration, two closed models; no intervals or
  repeats; some numbers disagree between text and figures (a plateau near
  66.5\% against the abstract's 68\%). \emph{Lines 787, 789.}
\item[Artifact.] No code released; a distilled set of about 31k examples is
  used to fine-tune small models, and the range does not say it is released.
  \emph{Lines 778, 789.}
\end{description}
A field that says the range is silent where it is silent (cost; variance
intervals; a released dataset) is correct. A field that contradicts a line
(for example an unlimited retry budget, or a test suite as the oracle) scores
2, quoting that line.

\end{document}
